\section{A Field That Needs the Other Service}
\label{sec:ch11-a-field-from-the-other-service}
\index{requires@\code{@requires}!on Mosaic|(}%

Mosaic sells furniture, lighting, kitchenware, textiles and storage, and it
ships them. A wardrobe goes on a pallet and a tea towel goes in a bag, so the
delivery charge follows from the category. Pricing owns what things cost.
Catalog owns what things are. Neither of them can work out the delivery charge
alone.

\index{computed field}%
Apollo's documentation calls a field like this a computed field: one worked out
from values of the same entity that another subgraph
resolves~\autocite{apollo2026contributefields}. This one goes a step beyond the
canonical example, because half its input is local. The category picks the
rate, and the price, which Pricing does own, decides whether the rate is
waived. A field that only echoed a remote value back would not need any of
this, because the router could read it from Catalog and never call Mosaic at
all.

Here is the whole of it, in \code{Pricing/Types/ProductPricingNode.cs}, beside
the \code{price} field that has been there since
chapter~\ref{ch:hotchocolate-quickly}:

\begin{minted}[fontsize=\footnotesize,breaklines]{csharp}
[Requires("category")]
public static async Task<Money> GetShippingCostAsync(
    [Parent] Product product,
    IPriceByProductIdDataLoader priceByProductId,
    CancellationToken cancellationToken)
{
    var category = product.Category
        ?? throw new InvalidOperationException(
            $"Product {product.Id} arrived without a category, so its shipping "
            + "cost cannot be worked out. The field requires 'category' and the "
            + "representation did not carry one.");

    var rate = ShippingRates.For(category);

    var price = await priceByProductId.LoadAsync(product.Id, cancellationToken);
    var amount = price?.Amount
        ?? throw new InvalidOperationException(
            $"No price is seeded for product {product.Id}.");

    return ShippingRates.IsWaived(category, amount.Amount)
        ? new Money(0m, amount.Currency)
        : new Money(rate, amount.Currency);
}
\end{minted}

\index{layout rule!fields live with their domain}%
It is in the Pricing folder, not the Catalog one, which is the layout rule this
repository has held since chapter~\ref{ch:hotchocolate-quickly} and the reason
chapter~\ref{ch:the-first-cut-extracting-catalog}'s extraction was a file move.
The rates are in \code{Pricing/Model/ShippingRates.cs}: 39.00 for furniture,
9.90 for lighting, 7.90 for storage, 5.90 for kitchen, 4.90 for textiles, and
free at or above 100 except for furniture, because a pallet costs what a pallet
costs. That threshold is a bare decimal and is compared against whatever the
price says, which is only safe because every price Mosaic holds is in euros.
Chapter~\ref{ch:hard-modeling-problems} is where a graph with two currencies in
it would have to do better than that.

The category is the other half, and it lives on the \code{Product} stub in the
\code{Catalog} folder, which is where a Catalog concept belongs even in a
service that does not own one:

\begin{minted}[fontsize=\footnotesize]{csharp}
[External]
[GraphQLNonNullType]
public ProductCategory? Category { get; set; }
\end{minted}

Three things in those three lines are worth stopping on, and two of them are
traps.

\index{external@\code{@external}!what it declares}%
\code{[External]} publishes the field and disclaims it in the same breath. The
type needs the field in its schema so that \code{@requires(fields: "category")}
has something to name, and the directive tells the composer that this is not a
place the value can be got from. Ask the router for \code{Product.category} and
it goes to Catalog. Nothing routes here for it, ever.

\index{external@\code{@external}!needs a setter}%
The property has a setter, and nothing in Mosaic ever assigns it. That is not
an oversight, it is the mechanism, and
section~\ref{sec:ch11-how-the-value-arrives} is about who does the assigning. A
get-only property here compiles, composes, and leaves the resolver holding
null for every product in the catalogue.

And the two nullabilities disagree on purpose. The schema says the field cannot
be null, because that is what Catalog says. The C\# property can be, because
most representations carry no category at all.
Section~\ref{sec:ch11-what-the-composer-insists-on} is where I found out what
happens if you let those two agree.

The enum itself is declared a second time, beside that stub, with the same five
members Catalog has. A copy, like \code{ProductKey} in
chapter~\ref{ch:the-first-cut-extracting-catalog} and for the same reason: the
values are a contract between two services, and a shared project would turn an
agreement about a wire format into a compile-time dependency. Unlike
\code{ProductKey} this copy is visible in the schema, so the composer has an
opinion about it, and chapter~\ref{ch:hard-modeling-problems} owns shared enums
as a subject.

\subsection*{What the schema says now}

\index{SDL!requires@\code{@requires} in the published schema}%
The published SDL, from \code{schema/mosaic.graphql} at tag \code{ch11}:

\begin{graphqlsdl}[fontsize=\footnotesize]
type Product @key(fields: "id") {
  availableQuantity: Int!
  price: Money!
  shippingCost: Money! @requires(fields: "category")
  reviews(first: Int, after: String, last: Int, before: String): ReviewConnection!
  averageRating: Float
  id: ID!
  category: ProductCategory! @external
}
\end{graphqlsdl}

The \code{reviews} arguments carry descriptions in the real file and are
collapsed here; nothing else is changed.

One line I did not write also changed. The \code{@link} at the top of the
schema grew two entries on its own, with nothing touched in \code{Program.cs}:

\begin{graphqlsdl}[fontsize=\footnotesize]
import: ["@external", "@key", "@requires", "@shareable", "@tag", "FieldSet"]
\end{graphqlsdl}

HotChocolate imports the directives it emitted. That is the same machinery
chapter~\ref{ch:composition} found putting \code{FieldSet} into the list, and
this time it is doing exactly what it should.

\subsection*{The plan stops being uniform}

\index{query plan!dependencies array@\code{dependencies} array!non-uniform}%
Chapter~\ref{ch:enter-the-router} read a query plan whose \code{dependencies}
array had four entries that all said the same thing, and closed by saying that
\code{@requires} is what should make one of them say something else. It does.
Ask the router for a title and a shipping cost, with no category anywhere in
the document, and the entity fetch carries this:

\begin{minted}[fontsize=\footnotesize,breaklines]{json}
{
  "coordinate": { "typeName": "Product", "fieldName": "shippingCost" },
  "isUserRequested": true,
  "dependsOn": [
    { "fetchId": 0, "subgraph": "catalog",
      "coordinate": { "typeName": "Product", "fieldName": "category" },
      "isKey": false, "isRequires": true },
    { "fetchId": 0, "subgraph": "catalog",
      "coordinate": { "typeName": "Product", "fieldName": "id" },
      "isKey": true, "isRequires": false }
  ]
}
\end{minted}

Two dependencies for one field, and the flags are how they differ:
\code{isKey} on the identifier, \code{isRequires} on the category. Every other
entry in the array still has the single key dependency
chapter~\ref{ch:enter-the-router} printed.

The \code{representations} array of that fetch grew a second entry too, which
is the same fact stated as a fragment the router will use to build each one:

\begin{minted}[fontsize=\footnotesize,breaklines]{json}
"representations": [
  { "kind": "@key", "typeName": "Product",
    "fragment": "fragment Key on Product { __typename id }" },
  { "kind": "@requires", "typeName": "Product", "fieldName": "shippingCost",
    "fragment": "fragment Requires_for_shippingCost on Product { category }" }
]
\end{minted}

\index{query plan!fetch count unchanged by requires@\code{@requires}}%
The number of fetches did not move. Two before, two after. The extra dependency
travels in a call that was already being made, which is
chapter~\ref{ch:enter-the-router}'s reading lesson holding up under the first
case that could have broken it: what costs you a round trip is a distinct
dependency that nothing else already satisfies, not a dependency as such.

\subsection*{What actually went over the wire}

\index{Advanced Request Tracing}%
Here I hit the problem chapter~\ref{ch:enter-the-router} left me with. Neither
Mosaic service logs request bodies, so I cannot do what
chapter~\ref{ch:how-a-federated-query-actually-runs} did and read the
representation out of a subgraph console. The router will tell me instead. Send
the query with \code{X-WG-Trace: true} and the response carries a \code{trace}
extension with every subgraph request in it, bodies included. It works for the
same reason the query plan headers work: development mode.

Two other spellings of that header, \code{X-WG-Trace-Enabled} and
\code{X-WG-Include-Trace}, are ignored without complaint, which cost me a few
minutes. Advanced Request Tracing as a subject belongs to
chapters~\ref{ch:inside-the-router-and-the-composer}
and~\ref{ch:observability}; here it is a tool.

The client asked for this, with no category in it:

\begin{minted}{graphql}
{ browseProducts(first: 2) { nodes { title shippingCost { amount } } } }
\end{minted}

What the router sent Catalog:

\begin{minted}[fontsize=\footnotesize,breaklines]{json}
{
  "query": "query($a: Int){browseProducts(first: $a){nodes {title category __typename id}}}",
  "variables": { "a": 2 }
}
\end{minted}

And what it then sent Mosaic:

\begin{minted}[fontsize=\footnotesize,breaklines]{json}
{
  "query": "query($representations: [_Any!]!){_entities(representations: $representations){... on Product {__typename shippingCost {amount}}}}",
  "variables": {
    "representations": [
      { "__typename": "Product", "category": "KITCHEN", "id": "UHJvZHVjdDoAAACgAAAAQIAAAAAAAAAW" },
      { "__typename": "Product", "category": "FURNITURE", "id": "UHJvZHVjdDoAAACgAAAAQIAAAAAAAAAE" }
    ]
  }
}
\end{minted}

\index{representation!key plus required fields}%
A representation is a key plus whatever some field required. That is the whole
mechanism, and the cost of it is visible in the first of those two bodies:
Catalog is being asked for a column nobody wants to see. The client never
receives \code{category} and has no way of knowing it was fetched.

\index{router!localhost fallback, seen from the trace}%
The same trace corroborates something chapter~\ref{ch:enter-the-router} could
only demonstrate by turning it off. Two entries under \code{load\_stats} for the
same fetch, with the timing fields every entry carries taken out:

\begin{minted}[fontsize=\footnotesize,breaklines]{text}
"get_conn":      { "host_port": "host.docker.internal:5101" }
"connect_done":  { "addr": "127.0.0.1:5101",
                   "err": "dial tcp 127.0.0.1:5101: connect: connection refused" }
\end{minted}

The router tried the URL as written, was refused, and went to the host instead.
That is the fallback setting narrating itself.

\subsection*{What it cost Mosaic}

\index{measurement!requires@\code{@requires} on the storefront}%
Mosaic still reports on itself, so this half is measurable. This is the
storefront query at twenty-five products, warm, with and without
\code{shippingCost} in the selection and nothing else different:

\begin{minted}[fontsize=\footnotesize,breaklines]{graphql}
{ browseProducts(first: 25) { nodes { title price { amount currency }
    shippingCost { amount currency } availableQuantity averageRating } } }
\end{minted}

\begin{minted}[fontsize=\footnotesize]{text}
                        resolvers   SQL
without shippingCost           76     3
with shippingCost             101     3
\end{minted}

\code{Product.reviews} is deliberately not in that selection, and leaving it
out is the difference between a number and a measurement. A resolver runs for
every review author, so a query that walks the reviews counts differently once
anything has written one, and this book's own Postman collection submits a
review every time it runs. \code{averageRating} still reaches the review table,
so nothing about the statement count is being avoided.

Twenty-five more resolvers, one per product, and no more database work at all.
The shipping resolver takes the same DataLoader that \code{price} takes, so the
second \code{LoadAsync} for a key finds it already resolved and the batch never
grows.

\index{gates!counts belong in one, timings do not}%
Both rows are steps in the verification script rather than numbers I typed out
of a terminal. Chapter~\ref{ch:enter-the-router} drew the line: a timing gets a
committed script a reader can run, because asserting milliseconds makes a gate
fail on a busier laptop, and a count gets an assertion, because a count is the
same on every machine. These are counts. The gate warms both documents, sends
each one once more, and reads the two timeline lines back out of Mosaic's own
console.

Catalog does slightly more work than it did, one more column in its projection,
and Catalog reports nothing. The cheapest half of this change is the half I can
measure, which is worth noticing about every number in this section.

\subsection*{Verify it in Postman}

\index{Postman!asserting a computed field}%
This chapter's collection is \code{mosaic-entities}: five requests and nineteen
assertions, two of the requests against the router and three against a subgraph
directly, because the subgraph is where you can ask a question the router will
not pass on. The rule rather than the rate is what it checks, so a change to
the table fails on the product it broke:

\begin{minted}[fontsize=\footnotesize,breaklines]{javascript}
pm.test("every product is priced for delivery by its category", function () {
    for (const node of nodes) {
        const rate = RATES[node.category];
        pm.expect(rate, `no rate for ${node.category}`).to.be.a("number");
        const waived = node.category !== "FURNITURE" && node.price.amount >= 100;
        pm.expect(node.shippingCost.amount, node.title).to.eql(waived ? 0 : rate);
        pm.expect(node.shippingCost.currency).to.eql(node.price.currency);
    }
});
\end{minted}

Two assertions beside it keep that one honest, and they are the sort of thing
easy to leave out: at least one product ships free, and not all of them do.
Without those, a rule that waived everything would pass. Of the twenty-five
seeded products, I counted nine that ship free, and the research note records
how to count them again.

\medskip
So the value arrives. What I have not said yet is how it gets from the
representation into that \code{product.Category} the resolver reads, and the
answer is not the one I expected.

\index{requires@\code{@requires}!on Mosaic|)}%
